\documentclass[9pt,a4paper]{extarticle}

% ======================
% Codificação, idioma e layout
% ======================
\usepackage[T1]{fontenc}
\usepackage[utf8]{inputenc}
\usepackage[brazil]{babel}
\usepackage{geometry}
\geometry{margin=2cm}
\usepackage{parskip}
\usepackage{setspace}
\usepackage{microtype}
\usepackage{enumitem}

% ======================
% Matemática e símbolos
% ======================
\usepackage{amsmath, amssymb, amsthm, bm}

% Algoritmos (CLRS-style)
\usepackage{algpseudocode}

% Cores
\usepackage{xcolor}

% ======================
% Figuras e tabelas
% ======================
\usepackage{graphicx}
\usepackage{booktabs}
\usepackage{array}
\usepackage{tabularx}
\usepackage{float}
\usepackage{caption}
\usepackage{subcaption}

% ======================
% Hiperlinks
% ======================
\usepackage{hyperref}
\hypersetup{
  colorlinks=true,
  linkcolor=blue,
  urlcolor=blue
}

% ======================
% Início do documento
% ======================
\begin{document}

\begin{center}
  \Large \textbf{Prova 2 \\ Algoritmos de Ordenção e Análise de algoritmo Iterativos e Recursivos} \\[4pt]
  \normalsize Disciplina: PCC104 - Projeto de Análise de Algoritmos
\end{center}

\vspace{1cm}

\textbf{Atenção: Você deverá entregar o código das implementações da lista de exercícios impresso juntamente com a prova. Algumas das questões farão referência ao seu próprio código. A não entrega do código implica em nota 0 nestas questões.}

\begin{enumerate}
  \item Baseado na sua implementação do \texttt{quicksort}, mostre a execução do método \texttt{partition} para:
  
  \[A = [2,8,7,1,3]\]
  \[p = 0\]
  \[r = 4\]

  \item Em quais linhas da sua implementação do método \texttt{partition} são feitas as trocas de elementos? Qual o custo das trocas?
  
  \item Construa um \texttt{max-heap} para o array $A = [4,1,3,2,16,9,10,14,8,7]$. Mostre o heap em forma de árvore e em forma de array.
  
  \item Considere o algoritmo abaixo, que verifica se o arranjo $A[1\,..\,n]$ contém algum par de elementos cuja soma seja zero.

    \begin{center}
    \begin{minipage}{0.65\textwidth}
    \textbf{\textsc{Soma-Zero}}$(A, n)$
    \begin{algorithmic}[1]
        \For{$i \gets 1$ \textbf{to} $n - 1$}
            \For{$j \gets i + 1$ \textbf{to} $n$}
                \If{$A[i] + A[j] = 0$}
                    \State \Return \textsc{verdadeiro}
                \EndIf
            \EndFor
        \EndFor
        \State \Return \textsc{falso}
    \end{algorithmic}
    \end{minipage}
    \end{center}

\begin{enumerate}[label=(\alph*)]
    \item Defina o tamanho da entrada do algoritmo.
    \item Determine o número de vezes que o teste do laço interno (linha 2) é executado no \textbf{pior caso}, em função de $n$. Descreva uma instância que produza o pior caso.

    \textit{Sugestão: para cada valor de $i$, conte quantos valores de $j$ são considerados e some sobre $i$. Use a identidade $\sum_{i=1}^{n-1} (n - i) = \dfrac{n(n-1)}{2}$.}

    \item Construa a tabela de custos no estilo CLRS para o pior caso e escreva $T_{\text{pior}}(n)$ como soma dos termos $c_i$ ponderados por suas frequências.
    \item Mostre que $T_{\text{pior}}(n) = \Theta(n^2)$, exibindo as constantes da definição assintótica.
    \item Descreva o \textbf{melhor caso} do algoritmo: quantas execuções da linha 3 ocorrem? Qual é a classe assintótica de $T_{\text{melhor}}(n)$?
\end{enumerate}

  \item Considere o algoritmo recursivo abaixo, que retorna a soma dos elementos do arranjo $A[p\,..\,r]$ pelo método de divisão e conquista.

    \begin{center}
    \begin{minipage}{0.7\textwidth}
    \textbf{\textsc{Soma}}$(A, p, r)$
    \begin{algorithmic}[1]
        \If{$p = r$}
            \State \Return $A[p]$
        \EndIf
        \State $q \gets \lfloor (p + r) / 2 \rfloor$
        \State $s_1 \gets \textsc{Soma}(A, p, q)$
        \State $s_2 \gets \textsc{Soma}(A, q + 1, r)$
        \State \Return $s_1 + s_2$
    \end{algorithmic}
    \end{minipage}
    \end{center}

    Considere $n = r - p + 1$ o tamanho do subarranjo na chamada corrente.

    \begin{enumerate}[label=(\alph*)]
        \item Identifique os passos de \textbf{dividir}, \textbf{conquistar} e \textbf{combinar} no algoritmo, no espírito do paradigma de divisão e conquista.
        \item Escreva a relação de recorrência $T(n)$, justificando o número de chamadas recursivas, o tamanho de cada subproblema e o custo de combinar os resultados.
        \item Resolva a recorrência.
        \item Mostre que $T(n) = \Theta(n)$.
    \end{enumerate}
\end{enumerate}


\end{document}

